\subsection{PRODUCT OF TWO INFINITE SUMS}

PROPOSITION 1. If the families \((x_{\lambda})_{\lambda \in L}\) and \((y_{\mu})_{\mu \in M}\) of finite real numbers are summable in \(\mathbf{R}\) , then so is the family \((x_{\lambda}y_{\mu})_{(\lambda ,\mu)\in L\times M}\) and we have

\[
\sum_ {(\lambda , \mu) \in \mathbf {L} \times \mathbf {M}} x _ {\lambda} y _ {\mu} = (\sum_ {\lambda \in \mathbf {L}} x _ {\lambda}) (\sum_ {\mu \in \mathbf {M}} y _ {\mu}).\tag{1}
\]

Every finite subset of \(\mathbf{L} \times \mathbf{M}\) is contained in a finite subset of the form \(\mathbf{H} \times \mathbf{K}\) , where \(\mathbf{H}\) is a finite subset of \(\mathbf{L}\) and \(\mathbf{K}\) is a finite subset of \(\mathbf{M}\) . By hypothesis, there exists a number \(a > 0\) such that \(\sum_{\lambda \in \mathbf{H}} |x_{\lambda}| \leqslant a\) and \(\sum_{\mu \in \mathbf{K}} |y_{\mu}| \leqslant a\) for all finite subsets \(\mathbf{H}\) and \(\mathbf{K}\) of \(\mathbf{L}\) and \(\mathbf{M}\) respectively; therefore

\[
\sum_ {(\lambda , \mu) \in \mathbf {H} \times \mathbf {K}} | x _ {\lambda} y _ {\mu} | = (\sum_ {\lambda \in \mathbf {H}} | x _ {\lambda} |) (\sum_ {\mu \in \mathbf {K}} | y _ {\mu} |) \leqslant a ^ {2},
\]

and this shows that the family \((x_{\lambda}y_{\mu})\) is summable in \(\mathbf{R}\) , by Theorems 1 and 3. By associativity we can write {[}Chapter III, § 5, no. 3, formula (2){]}

\[
\sum_ {(\lambda , \mu) \in \mathbf {L} \times \mathbf {M}} x _ {\lambda} y _ {\mu} = \sum_ {\lambda \in \mathbf {L}} (\sum_ {\mu \in \mathbf {M}} x _ {\lambda} y _ {\mu}) = \sum_ {\lambda \in \mathbf {L}} x _ {\lambda} (\sum_ {\mu \in \mathbf {M}} y _ {\mu}) = (\sum_ {\lambda \in \mathbf {L}} x _ {\lambda}) (\sum_ {\mu \in \mathbf {M}} y _ {\mu});
\]

hence the result.

